\section{Methodik}
\label{sec:methodik}

\subsection{Werkzeugkette}
\label{sec:toolchain}

Die Kette besteht aus vier Stufen: Bearbeiten, Kompilieren, Auswerten und
Veröffentlichen. Abbildung~\ref{fig:pipeline} stellt sie dar.

\begin{figure}[htbp]
  \centering
  \begin{tikzpicture}[
      >=Stealth,
      stage/.style={draw, rounded corners=3pt, minimum width=26mm, minimum height=10mm,
                    align=center, fill=blue!8},
      link/.style={->, thick, black!65},
    ]
    \node[stage] (edit) {Bearbeiten};
    \node[stage, right=14mm of edit] (build) {Kompilieren};
    \node[stage, right=14mm of build] (check) {Auswerten};
    \node[stage, right=14mm of check] (pub) {Veröffentlichen};
    \draw[link] (edit) -- (build);
    \draw[link] (build) -- (check);
    \draw[link] (check) -- (pub);
    \draw[link, red!70!black, dashed] (check.south) -- node[below=2mm, font=\small]
      {Fehler} (build.south);
  \end{tikzpicture}
  \caption{Vierstufige Dokumentenpipeline mit Rückkopplung.}
  \label{fig:pipeline}
\end{figure}

\begin{warnung}{Hinweis zu shell-escape}
  Makros, die externe Programme starten, benötigen
  \texttt{shell-escape}. Diese Option ist standardmäßig \emph{ausgeschaltet} und
  kann in den Einstellungen von \emph{LaTeX Studio} aktiviert werden.
\end{warnung}

\subsection{Messdesign}

Für jede Stufe wird die benötigte Zeit in Sekunden gemessen. Der Mittelwert
berechnet sich als
\[
  \bar{t} \;=\; \frac{1}{n}\sum_{i=1}^{n} t_i,
\]
die Standardabweichung entsprechend nach der üblichen Formel. Die Stichprobe
umfasst $n = \num{30}$ Läufe pro Stufe.

\subsection{Datenfluss}

\begin{table}[htbp]
  \centering
  \caption{Eingangsdaten der Auswertung.}
  \label{tab:inputs}
  \begin{tabular}{l l r}
    \toprule
    Quelle & Inhalt & Umfang \\
    \midrule
    \texttt{main.log}     & Meldungen des Satzsystems & \SI{180}{\kibi\byte} \\
    \texttt{main.aux}     & Labels und Verweise      & \SI{12}{\kibi\byte} \\
    \texttt{references.bib} & Literaturreferenzen      & \SI{4}{\kibi\byte} \\
    \bottomrule
  \end{tabular}
\end{table}
